\section{Bối cảnh}

Kiến trúc cloud thường tách workload thành nhiều VPC để phân chia miền lỗi,
quyền quản trị và chính sách bảo mật. Khi các workload cần trao đổi bằng địa
chỉ riêng, hệ thống phải bổ sung kết nối liên VPC. Hai phương án AWS được khảo
sát trong đồ án là VPC Peering và Transit Gateway (TGW). Peering tạo kết nối
trực tiếp giữa từng cặp VPC, trong khi TGW cung cấp một hub định tuyến tập
trung cho nhiều attachment \cite{aws-vpc-peering,aws-tgw-overview}.

Luồng vấn đề của nghiên cứu được tóm tắt như sau:

\begin{center}
Nhiều VPC $\rightarrow$ cần kết nối liên VPC $\rightarrow$
Peering hoặc TGW $\rightarrow$ kiến trúc đường đi khác nhau
$\rightarrow$ có thể tạo khác biệt hiệu năng.
\end{center}

Khác biệt kiến trúc không tự động chứng minh khác biệt hiệu năng. Phép so sánh
cần giữ nguyên endpoint, vị trí triển khai, tải và công cụ đo; chỉ thay đường
định tuyến giữa hai mode.

\section{Vấn đề nghiên cứu}

Vấn đề trung tâm là: với cùng endpoint EC2 và cùng điều kiện benchmark, hiệu
năng mạng quan sát được giữa VPC Peering và TGW khác nhau như thế nào? Nghiên
cứu không đo hai môi trường độc lập vì khác biệt instance có thể làm nhiễu kết
quả. Thay vào đó, Terraform chuyển route của cùng bốn EC2 giữa hai đường kết
nối rồi kiểm tra lại trước khi đo.

\section{Mục tiêu}

Đồ án đánh giá các đại lượng sau:

\begin{itemize}
  \item RTT khi mạng rảnh giữa A và B;
  \item receiver goodput của TCP một luồng;
  \item số TCP retransmission do iperf3 ghi nhận;
  \item tổng receiver goodput xấp xỉ trong kịch bản fan-in;
  \item CPU busy tại endpoint gửi và nhận;
  \item thay đổi các network allowance counter do ENA cung cấp.
\end{itemize}

Mục tiêu phân tích không chỉ là xếp hạng hai giá trị throughput. Báo cáo còn
phải xác định khi nào endpoint trở thành bottleneck và ranh giới nào khiến dữ
liệu chưa đủ để suy ra năng lực tối đa của dịch vụ AWS.

\section{Câu hỏi nghiên cứu}

\begin{enumerate}
  \item RTT qua Peering và TGW khác nhau bao nhiêu trong cấu hình thử nghiệm?
  \item Receiver goodput TCP một luồng khác nhau như thế nào?
  \item Khi A, C và D cùng gửi tới B, tổng goodput thay đổi như thế nào?
  \item Endpoint EC2/ENA ảnh hưởng ra sao tới cách diễn giải kết quả fan-in?
  \item Trade-off giữa performance và khả năng mở rộng kiến trúc là gì?
\end{enumerate}

\section{Phạm vi}

Hệ thống thực nghiệm đặt tại region \texttt{us-east-1}, Availability Zone
\texttt{us-east-1a}; gồm bốn VPC và bốn EC2 \texttt{c6i.large}. Traffic đo đi
qua private IP, dùng TCP và MTU 1500 byte cho baseline. RTT dùng ICMP. Mỗi bộ
đo chính có năm lượt.

Phạm vi này không đại diện cho mọi region, Availability Zone, loại instance,
MTU, số luồng hay quy mô VPC. Kết quả không phải benchmark năng lực tối đa của
Peering hoặc TGW. Chi phí được thảo luận theo mô hình tính phí, không dùng như
số liệu hóa đơn của phiên đo.

\section{Đóng góp của đồ án}

\begin{itemize}
  \item Hạ tầng bốn VPC có thể tái tạo bằng Terraform.
  \item Cùng endpoint được giữ nguyên khi thay đổi đường Peering và TGW.
  \item Ba kịch bản chính bao phủ RTT, TCP single-flow và TCP fan-in.
  \item Raw data, metadata, route snapshot, CPU và ENA counter được lưu cùng
        kết quả.
  \item Script độc lập tính lại số liệu và kiểm tra tính nhất quán của nguồn.
  \item Phân tích tách kết quả thực nghiệm khỏi đặc tính kiến trúc và nêu rõ
        ảnh hưởng của bottleneck tại máy nhận.
\end{itemize}

\section{Cấu trúc báo cáo}

Chương 2 trình bày nội dung kỹ thuật trực tiếp phục vụ phép đo. Chương 3 mô tả
hạ tầng, thiết kế đối chứng, quy trình benchmark và kết quả thực nghiệm chi
tiết. Chương 4 tổng hợp các nghiên cứu và tài liệu liên quan. Chương 5 đánh giá
các finding, trả lời câu hỏi nghiên cứu, nêu hạn chế, hướng phát triển và kết
luận.
